\documentclass[12pt]{article}
\usepackage{amsmath,amssymb,booktabs,longtable}
\usepackage{fontspec}
\usepackage{polyglossia}
\setmainlanguage{arabic}
\setotherlanguage{english}
\newfontfamily\arabicfont[Script=Arabic]{Amiri}
\usepackage{geometry}
\geometry{margin=2.2cm}
\usepackage{hyperref}
\usepackage{svg}

\title{مقارنة عجلات الأساس 29 و97 و251 و997 في نموذج الارتباط الزوجي النقطي}
\author{مشروع تجربة البحث عن الأعداد الأولية}
\date{2026-08-31}

\begin{document}
\maketitle

\section{السؤال}

نكرر نموذج الغربلة العشوائية النقطي نفسه عند عجلات أساس
\[
q_0\in\{29,97,251,997\},
\]
ثم نواصل الغربلة العشوائية حتى حد \(Q\)، ونقيس أين نحصل على أفضل توازن بين:

\begin{enumerate}
\item مقدار التشتت داخل الفترات؛
\item الارتباط الذاتي عند الإزاحة \(1\) بين الفترات المتجاورة.
\end{enumerate}

العينة ثابتة في جميع التجارب: 450 فترة
\[
p_n^2 < x < p_{n+1}^2
\]
مع \(p_n\ge 4001\).

\section{النموذج}

للعجلة الثابتة \(q_0\)، نرمز إلى عدد المواقع الناجية في الفترة \(i\) بـ
\[
C_i(q_0).
\]

ويظل متوسط عدد الأوليات المستهدف هو المتوسط السلس المعاير عالميًا:
\[
\mu_i.
\]

نستخدم تباين عجلة الأساس:
\[
V_i(q_0)
=
\mu_i\left(1-\frac{\mu_i}{C_i(q_0)}\right).
\]

ثم نعرّف الباقي المعياري:
\[
Z_i
=
\frac{X_i-\mu_i}{\sqrt{V_i(q_0)}}.
\]

ومقياس التشتت:
\[
R
=
\frac{1}{m-1}\sum_{i=1}^{m} Z_i^2.
\]

بعد عجلة الأساس نختار، لكل أولي
\[
q_0<r\le Q,
\]
فئة ممنوعة واحدة modulo \(r\) عشوائيًا، كما في تجربة ذيل
Hardy--Littlewood السابقة.

إذا كان حاصل بقاء الذيل المتوقع هو
\[
S(q_0,Q)
=
\prod_{q_0<r\le Q}\left(1-\frac1r\right),
\]
فإننا نطبق thinning داخل كل فترة باحتمال
\[
\theta_i
=
\frac{\mu_i}{C_i(q_0)S(q_0,Q)}
\]
حتى يبقى المتوسط المتوقع مساويًا لـ \(\mu_i\).
وبذلك نحاول عزل أثر البنية الزوجية بدل تغيير نموذج المتوسط.

\section{مقياس التوازن}

لكل إعداد نحسب متوسطات Monte Carlo وانحرافاتها المعيارية.
ثم:
\[
z_R
=
\frac{\overline R_{\mathrm{MC}}-R_{\mathrm{obs}}}
{\sigma_R},
\qquad
z_\rho
=
\frac{\overline \rho_{\mathrm{MC}}-\rho_{\mathrm{obs}}}
{\sigma_\rho}.
\]

ومقياس التوازن الأساسي:
\[
B
=
\sqrt{z_R^2+z_\rho^2}.
\]

القيمة الأصغر أفضل. هذا المقياس ليس ``نسبة عشوائية''، بل مسافة معيارية
بين النموذج والبيانات في بعدين محددين.

\section{المسح الأولي}

أُجري مسح على عدة قيم لـ \(Q\) حتى \(7919\)، مع 400 محاكاة نهائية تقريبًا
لكل مرحلة. أعطى المسح منطقة منافسة واضحة للعجلات \(29,97,251\)،
بينما ظلت عجلة \(997\) ضعيفة خصوصًا في تفسير الارتباط السلبي.

بسبب ضوضاء Monte Carlo أُجريت جولات تأكيد مستقلة للإعدادات الأقوى.

\section{نتائج التأكيد}

\begin{center}
\begin{tabular}{rrrrrrrr}
\toprule
\(q_0\) & \(Q\) & \(R_{\rm obs}\) & \(R_{\rm MC}\) &
\(\rho_{\rm obs}\) & \(\rho_{\rm MC}\) & \(B\) & ملاحظة\\
\midrule
251 & 7919 & 0.666831 & 0.646079 & -0.123707 & -0.068098 & 1.298 & أفضل \(B\)\\
251 & 6311 & 0.666831 & 0.659195 & -0.123707 & -0.059091 & 1.440 & قريب\\
29  & 7919 & 0.447870 & 0.479401 & -0.123661 & -0.074912 & 1.448 & أفضل ارتباط\\
97  & 7919 & 0.546196 & 0.572568 & -0.123669 & -0.065409 & 1.509 & قريب\\
997 & 1259 & 0.989424 & 0.967225 & -0.124267 & -0.007610 & 2.553 & ضعيف زوجيًا\\
\bottomrule
\end{tabular}
\end{center}

\section{قيم الاختبار التجريبية}

في جولات التأكيد:

\begin{itemize}
\item \(q_0=251,Q=7919\):
\[
p_R=0.652,\qquad p_\rho=0.238,\qquad p_{\rm runs}=0.054.
\]

\item \(q_0=29,Q=7919\):
\[
p_R=0.310,\qquad p_\rho=0.326,\qquad p_{\rm runs}=0.057.
\]

\item \(q_0=97,Q=7919\):
\[
p_R=0.483,\qquad p_\rho=0.179,\qquad p_{\rm runs}=0.046.
\]

\item \(q_0=997,Q=1259\):
\[
p_R=0.722,\qquad p_\rho=0.0075,\qquad p_{\rm runs}=0.0092.
\]
\end{itemize}

\section{ما هو ``الأفضل''؟}

إذا استعملنا مقياس المسافة المعيارية \(B\) المحدد مسبقًا، فإن أفضل إعداد
من الإعدادات المؤكدة هو:
\[
\boxed{q_0=251,\qquad Q=7919.}
\]

لكن الفرق بين \(q_0=251\) و\(q_0=29\) ليس كبيرًا.
إذا استعملنا بدلًا من ذلك معيار ``تعظيم أسوأ قيمة \(p\) بين \(R\) و\(\rho\)''
فإن \(q_0=29,Q=7919\) يصبح جذابًا أكثر.

لذلك الاستنتاج الأقوى ليس أن \(251\) رقم سحري، بل:
\[
\boxed{\text{أفضل منطقة في هذه التجربة تقع بين عجلة أساس صغيرة إلى متوسطة، وليست عند }997.}
\]

\section{المعنى البنيوي}

عجلة \(997\) تثبّت مسبقًا جزءًا كبيرًا من البنية الحسابية التي نريد من
ذيل singular series أن يفسرها. ولذلك يستطيع نموذجها مطابقة التشتت بسهولة،
لكنه يبقى قريبًا من الاستقلال في الارتباط:
\[
\rho_{\rm MC}\approx -0.008
\]
مقابل
\[
\rho_{\rm obs}\approx -0.124.
\]

عند خفض عجلة الأساس، يبقى قدر أكبر من البنية داخل الجزء العشوائي الزوجي.
لذلك يتحرك الارتباط بوضوح في الاتجاه الصحيح:
\[
-0.008\longrightarrow -0.065,\,-0.068,\,-0.075
\]
بحسب عجلة الأساس والإعداد.

لكن حتى أفضل النماذج لا تصل إلى الارتباط المرصود كاملًا.

\section{اختبار التتابعات يكشف بقايا بنية}

حتى في أفضل الإعدادات بقي عدد التتابعات في المحاكاة قرب
\[
233\text{--}236
\]
بينما المرصود:
\[
256.
\]

وقيم \(p\) كانت قرب الحد \(0.05\).

هذا مهم لأن مطابقة \(R\) والارتباط عند الإزاحة \(1\) لا تكفي لضمان مطابقة
البنية الكاملة للسلسلة. ما زال هناك أثر في ترتيب الإشارات لا يفسره النموذج
الزوجي الحالي تمامًا.

\section{الخلاصة}

التجربة تدعم ثلاث نتائج:

\begin{enumerate}
\item تثبيت عجلة كبيرة مثل \(997\) قبل إدخال pair model يخفي جزءًا مهمًا من البنية التي نريد قياسها.
\item عجلات الأساس \(29,97,251\) تعطي توازنًا أفضل بكثير.
\item أفضل إعداد وفق مقياس \(B\) الحالي هو \(251\to7919\)، لكن \(29\to7919\) منافس قوي، ولذلك لا يوجد دليل على أن \(251\) حد مميز جوهريًا.
\end{enumerate}

السؤال التالي الطبيعي هو توسيع مسح عجلة الأساس نفسها إلى:
\[
2,3,5,7,11,17,29,47,71,97,127,173,251,\ldots
\]
مع تثبيت \(Q=7919\)، لرسم منحنى ``مقدار البنية الذي نجمّده في العجلة''
مقابل ``مقدار البنية الذي يترك للـ pair model''.

\section{ملفات التجربة}

\begin{itemize}
\item \texttt{base\_wheel\_pair\_balance\_sweep.csv}
\item \texttt{base\_wheel\_pair\_balance\_best\_per\_base.csv}
\item \texttt{base\_wheel\_pair\_balance\_confirmation.csv}
\item \texttt{analyze\_base\_wheel\_pair\_balance.py}
\item \texttt{plots/base\_wheel\_pair\_balance\_score.svg}
\item \texttt{plots/base\_wheel\_pair\_balance\_confirmation.svg}
\end{itemize}

\end{document}
